% Capitulo I: Introduccion (maximo 3 paginas)

\chapter{INTRODUCCIÓN}

\section{Presentación del tema de proyecto}

El presente proyecto se ubica en el sector de los sistemas de aire comprimido de uso industrial. El aire comprimido es un servicio auxiliar presente en plantas de manufactura, minería, alimentos y otros rubros, donde se emplea para accionar válvulas, actuadores y herramientas neumáticas, así como en procesos de soplado y transporte. Un sistema de aire comprimido está conformado por equipos de generación, tratamiento y distribución, y su elemento principal es el compresor. En aplicaciones de potencia media son de uso extendido los compresores de tornillo rotativo lubricados de velocidad fija, que regulan la entrega de aire alternando entre los estados de carga y descarga dentro de una banda de presión \cite{liang2019}. Estos equipos suelen operar de forma continua y cuentan con un controlador electrónico que supervisa variables como presiones, temperaturas, horas de funcionamiento y alarmas. En ese contexto, la oportunidad que aborda el proyecto se ubica en la evaluación de la capacidad y del desempeño de estos compresores frente a la demanda de aire de la planta.

\section{Descripción de la situación problemática}

En una planta industrial, el compresor debe entregar el caudal que demandan los equipos a una presión suficiente. Cuando la demanda supera su capacidad, la presión cae por debajo de la mínima de operación y los equipos neumáticos se detienen o trabajan mal. La demanda, sin embargo, no es fija: las plantas incorporan nuevas máquinas, amplían turnos o acumulan fugas, y el compresor que antes bastaba puede dejar de hacerlo.

Para conocer el estado de la sala de compresores se realizan auditorías energéticas. En ellas se registran la corriente y la presión durante varios días con registradores portátiles, y con los datos del fabricante se estiman el caudal entregado, las horas en carga y descarga y la energía consumida; con ese perfil se evalúa si el sistema es eficiente y si soporta una demanda adicional \cite{cac2026}. La auditoría es puntual: exige instalar y retirar equipos en planta, y su resultado describe solo la semana medida. Si después cambia la demanda o el compresor pierde caudal por desgaste, la evaluación queda desactualizada y la decisión de aumentar la demanda, instalar un tanque o adquirir otro compresor se toma sin información vigente.

El problema afecta al cliente, que necesita saber si su sala soporta los nuevos equipos antes de instalarlos; al área de mantenimiento, que no advierte a tiempo una pérdida de capacidad; y al servicio técnico, que debe repetir visitas para actualizar cada evaluación. Ocurre en plantas de operación continua, con demandas que varían entre turnos y días, y personal que no siempre cuenta con formación en análisis de datos.

Los compresores actuales ya registran en su controlador el estado de carga, la presión, las temperaturas y las horas de operación, y pueden entregar esos datos por comunicación industrial \cite{pinzaapolo2026}. Se requiere, por tanto, una solución técnica que use estos datos para mantener actualizada la evaluación del compresor y que permita probar cambios antes de hacerlos en la planta.

\section{Formulación del proyecto}

Las evaluaciones de capacidad de los compresores de tornillo se hacen de forma puntual y pierden vigencia cuando cambia la demanda o el estado del equipo. Por ello, el problema de ingeniería se formula así: ¿cómo diseñar un sistema que evalúe continuamente la capacidad de un compresor de tornillo ante cambios de demanda?

\section{Objetivo general del proyecto}

Desarrollar una sombra digital para evaluar la capacidad y el desempeño operacional de un compresor de tornillo lubricado de velocidad fija.

\section{Alcances del proyecto}

El proyecto llega hasta la fase de diseño conceptual, con una simulación en MATLAB que verifica su funcionamiento en una planta virtual construida con los datos de una auditoría energética. La solución se plantea para un compresor de tornillo lubricado de velocidad fija, con regulación por carga y descarga y controlador electrónico con comunicación Ethernet. Los entregables son los requerimientos; la matriz morfológica y la evaluación de las propuestas; la arquitectura por niveles; el modelo de la sombra digital, que se calibra con siete días de datos; y la simulación de sus tres funciones: estado actual, capacidad ante una demanda adicional o un tanque mayor, y alerta de condición. Como caso de estudio se emplea un compresor de 29,2\,cfm y 9,5\,kW en carga, cuyos datos provienen de una auditoría energética de siete días, anonimizada. No se incluye la conexión a equipos reales; los compresores de velocidad variable, las salas con varios compresores y el control automático del equipo quedan como trabajo futuro.
